\chapter{Formal Semantics, Types, and Contracts}
\label{ch:semantics}

Here is a question you cannot currently answer about your own agent.

Given its configuration, will it terminate? The question is about the
configuration itself, rather than about how it behaved on the cases you tried. Most teams answer with a step limit, which is an
admission that the question was never asked. A step limit does not tell you
whether the agent halts. It tells you when to stop waiting.

Chapter~\ref{ch:terms} defined the vocabulary operationally, saying what a system
must track and enforce. That is enough to build with and not enough to prove
anything. This chapter covers the formal side, which arrived late to agent work
and arrived quickly. Over the past year the field acquired typed calculi for
agent composition, process-calculus semantics for tool protocols, information-flow
calculi for injection resistance, and contract frameworks that make resource
bounds part of a specification rather than a runtime hope.

The practical payoff is narrower than the theory suggests and still worth having.
You get lint rules that catch structural errors before deployment, a precise
statement of what a protocol translation loses, and a way to say what an agent
is allowed to consume that a scheduler can check.

\section{Why Agent Frameworks Resist Formalization}

An ordinary program has a semantics because every construct has defined behavior.
Agent frameworks break this in four places.

\textbf{Oracle calls.} A model invocation returns an unpredictable value from an
enormous space. Its behavior is not a function of its inputs in any usable sense.

\textbf{Unbounded loops.} The ReAct loop runs until the agent decides to stop,
and the agent deciding is the same oracle. Termination is not structural.

\textbf{Probabilistic choice.} Sampling makes two runs of the same configuration
different programs in the operational sense.

\textbf{Mutable shared environment.} Tools read and write state outside the
program, concurrently with other agents.

Any calculus for agents has to model all four rather than assume them away.
That is the bar the recent work clears.

\section{A Typed Calculus for Agent Composition}

The most direct result extends the simply-typed lambda calculus with exactly
those four features. The $\lambda_A$ calculus adds oracle calls, \emph{bounded fixpoints} to model the ReAct loop, probabilistic choice, and mutable environments \cite{lambdaA2026}. A bounded fixpoint is a recursive definition carrying an explicit iteration limit, which is how the calculus represents a loop that must stop. On this it proves \emph{type safety}, meaning a well-typed configuration cannot reach a stuck state, together with termination of those bounded loops.

The development is mechanized in Coq, a proof assistant that checks each step by machine. That matters here because hand proofs about probabilistic effects and mutable state are exactly the kind that look right and are not.

Two consequences are worth separating.

The theoretical one is that agent composition admits a semantics at all. Given a
configuration, you can ask whether it is well-formed and whether its loops
terminate, and get an answer that does not depend on running it.

The practical one is that lint rules fall out of the operational semantics rather
than out of taste. A rule derived this way comes with a soundness argument, so a
warning means something specific. Applied to a corpus of real agent
configurations, this catches structural errors that no amount of prompt review
would surface, because they are properties of the wiring rather than of the text.

The limitation is the bounded fixpoint. Termination is proved \emph{because} the
loop is bounded, which returns us to the step limit with a better justification.
What the calculus adds is that the bound is now part of the type, so a
configuration that omits it fails to typecheck instead of running until someone
notices.

\section{Process Calculus for Tool Protocols}

Tool protocols are where agents meet the outside world, and the field has
converged on a small number of them without anyone checking whether they express
the same things.

A \emph{process calculus} is a formal language for describing concurrent systems that communicate over channels, with rules for when two such systems count as equivalent. Applying one to two dominant approaches, Schema-Guided Dialogue and the Model Context Protocol, proves them structurally \emph{bisimilar} under an explicit mapping \cite{toolcalculus2026}. Bisimilar means that for practical
purposes each can simulate the other, so a system built on one is not
architecturally committed in the way it might appear.

The interesting half is the reverse direction. The inverse mapping is partial and
lossy, which means the translation back discards structure. If you are building a
bridge between protocols, the formalization tells you precisely what falls
through, and it is better to learn that from a proof than from an incident.

A related line formalizes agent communication itself under resource constraints,
producing a calculus that keeps the semantic richness of older agent
communication languages while remaining tractable for constrained deployment
\cite{uacp2026}. The recurring theme is that expressiveness and cost are formally
related, which is the same trade Chapter~\ref{ch:cost} makes empirically.

\section{Information Flow and Injection Resistance}

The most consequential formal result for practitioners concerns prompt injection.

Chapter~\ref{ch:capability} argues that the defense against injection is
structural, keeping trusted and untrusted content apart so that data cannot
become instruction. That argument can now be made precisely. \emph{Information-flow control} is the discipline of labelling data by origin and restricting how labels may propagate, so that content from an untrusted source cannot reach a position where it is obeyed. The LLMbda calculus is an untyped call-by-value lambda calculus that models agents, conversations, and information flow on that basis, giving a formal account of provenance-based defenses including the dual-LLM pattern \cite{llmbda2026}.

The value is in what it exposes about existing defenses. Flow tracking is easy to
implement incorrectly, deliberate relaxations are hard to audit, and the dual-LLM
pattern is usually hard-wired into an architecture rather than derived. A
calculus makes each of those a statement you can check. A relaxation becomes an
explicit rule in the system rather than a comment in a code review.

If you take one thing from this section, take the reframing. Injection resistance
is an information-flow property, and information flow is a well-studied area with
sixty years of results. Treating it as a prompt-engineering problem discards all
of them.

\section{Contracts and Resource-Bounded Specification}

Traditional software specifies behavior with contracts, meaning preconditions,
postconditions, invariants, and types. Agents typically specify behavior with a
paragraph of English in a system prompt. The gap between those two is where
governance failures come from.

Two frameworks close it from different directions.

\textbf{Behavioral contracts.} An agent behavioral contract is a tuple of
preconditions, invariants, governance policies, and recovery mechanisms, all
runtime-enforceable rather than documentary \cite{agentcontracts2026}. Because
agent behavior is stochastic, satisfaction is defined probabilistically with
explicit tolerance parameters, which is the honest formulation. A contract that
demanded deterministic satisfaction would be unsatisfiable by construction and
therefore useless.

\textbf{Resource contracts.} The older Contract Net Protocol coordinated
multi-agent work through contracts for task allocation. Modern agent protocols
standardize connectivity and say nothing about consumption. Agent Contracts
extend the metaphor to resource-bounded execution, unifying input and output
specifications, multi-dimensional resource constraints, temporal boundaries, and
success criteria under one lifecycle \cite{resourcecontracts2026}.

That second framework connects directly to Chapter~\ref{ch:budgets}. A budget
enforced by the policy is a good design. A budget expressed as a contract is a
budget a scheduler, a marketplace, or a counterparty can also check, which is
what you need once agents consume resources on behalf of parties who are not in
the room.

Formal verification has begun reaching the packaging layer as well. Work on agent
skill manifests defines a verification lattice running from unverified through
declared and tested to mechanically checked, and supplies methods for the top
level \cite{skillverify2026}. Given the supply-chain exposure described in
Chapter~\ref{ch:capability}, a mechanically checkable manifest is worth more than
a signature from a publisher nobody has audited.

\section{What Formalization Buys, and What It Does Not}

Be clear-eyed about the return.

\textbf{It buys structural error detection.} Configuration errors that are
properties of the wiring get caught statically. This is real, cheap, and
underused.

\textbf{It buys precise interface claims.} Bisimilarity results tell you when two
protocols are interchangeable and exactly what a translation loses.

\textbf{It buys auditable relaxations.} An information-flow calculus turns
``we allow this one exception'' into a rule with a name.

\textbf{It buys checkable resource claims.} A resource contract is enforceable by
a party other than the agent.

\textbf{It does not buy correctness of model outputs.} No calculus makes an oracle
call return the right answer. Every result here treats the model as an oracle and
reasons about the structure around it, which is the same posture the rest of this
book takes.

\textbf{It does not remove the need for runtime enforcement.} Type safety is a
property of configurations, not of adversaries. Chapter~\ref{ch:gates} and
Chapter~\ref{ch:alignment} still apply in full.

The honest summary is that formalization moves a class of failures from runtime
to compile time, and that class is smaller than the theory papers imply and
larger than most teams currently catch.

\section{Fallacies and Pitfalls}

\textbf{Fallacy. A step limit proves termination.} It bounds waiting. A bounded
fixpoint in a typed calculus proves termination because the bound is part of the
type, and a configuration missing it fails to typecheck \cite{lambdaA2026}.

\textbf{Fallacy. Two protocols that look similar are interchangeable.} They may
be bisimilar in one direction and lossy in the other. The direction you are
translating matters \cite{toolcalculus2026}.

\textbf{Fallacy. Injection is a prompting problem.} It is an information-flow
problem, and treating it otherwise discards a mature literature
\cite{llmbda2026}.

\textbf{Fallacy. A contract must be satisfied deterministically.} Agent behavior
is stochastic, so useful contracts carry explicit probabilistic tolerances
\cite{agentcontracts2026}.

\textbf{Pitfall. Formalizing the model rather than the structure.} The oracle is
not the tractable part. Every result in this chapter formalizes what surrounds
the oracle.

\textbf{Pitfall. Treating mechanization as optional.} Proofs about probabilistic
effects and mutable state are exactly where hand arguments fail. Prefer results
that carry machine-checked developments.

\section{Summary}

Agent frameworks resisted formalization for four specific reasons, namely oracle
calls, unbounded loops, probabilistic choice, and mutable shared environments.
Recent calculi model all four rather than assuming them away.

Typed composition gives well-formedness and termination as checkable properties,
with lint rules derived from the operational semantics instead of from
convention. Process-calculus treatment of tool protocols establishes when two
protocols are interchangeable and what a translation between them loses.
Information-flow calculi restate injection resistance as a provenance property,
which connects it to a mature body of work and makes deliberate relaxations
auditable.

Contracts supply the specification layer that agent frameworks lack. Behavioral
contracts carry preconditions, invariants, governance, and recovery with
probabilistic satisfaction. Resource contracts make consumption bounds
externally checkable, which is what Chapter~\ref{ch:budgets} needs once budgets
cross an organizational boundary.

None of this makes a model correct. All of it makes the structure around the
model something you can reason about before it runs.

\section*{Bibliographic Notes}

The typed calculus for agent composition, including the Coq development and the
derived lint tool, is \cite{lambdaA2026}. Process-calculus semantics for tool
protocols, with the bisimilarity result between Schema-Guided Dialogue and the
Model Context Protocol, is \cite{toolcalculus2026}. Resource-constrained agent
communication is formalized in \cite{uacp2026}.

The LLMbda calculus \cite{llmbda2026} builds on a long tradition of
information-flow control; readers who want the background should start with the classical treatment of provenance and of noninterference, the property that untrusted inputs cannot influence trusted outputs, before reading the agent adaptation.

Behavioral contracts are introduced in \cite{agentcontracts2026}, and
resource-bounded agent contracts in \cite{resourcecontracts2026}, which
explicitly positions itself against the Contract Net Protocol of 1980. Formal
verification of agent skill manifests appears in \cite{skillverify2026}.

Readers looking for the limits of runtime formal methods should read
Chapter~\ref{ch:alignment}, which covers a coverage bound on monitor recall that
constrains what any fixed-invariant runtime monitor can achieve.

\section*{Exercises}

\begin{enumerate}
\item Write down your agent's main loop as a bounded fixpoint. State the bound
explicitly and say what happens at the bound. If your answer is ``it returns
whatever it has,'' say what type that value has and whether callers handle it.

\item Take two tool protocols your system speaks. Construct the mapping from one
to the other and identify one piece of structure the reverse mapping loses.
Compare your answer to the lossiness result in \cite{toolcalculus2026}.

\item Express one injection defense in your system as an information-flow rule
over provenance labels. Identify every deliberate relaxation and say who
approved it.

\item Write a behavioral contract for one agent, giving preconditions,
invariants, governance policies, and recovery. Choose the probabilistic tolerance
and justify the number from the cost of violation rather than from convention.

\item Your agent consumes budget on behalf of a customer. Write the resource
contract that a third party could check without trusting the agent, following
\cite{resourcecontracts2026}. State which terms a scheduler can enforce and which
require post-hoc audit.

\item Argue for or against the claim that a mechanically checked skill manifest
is worth more than a publisher signature, using the supply-chain threat model of
Chapter~\ref{ch:capability}.
\end{enumerate}

\section*{Bibliography}

\begin{thebibliography}{20}

\bibitem{agentcontracts2026}
Varun Pratap Bhardwaj.
``Agent Behavioral Contracts: Formal Specification and Runtime Enforcement for Reliable Autonomous AI Agents.''
arXiv:2602.22302, 2026.

\bibitem{lambdaA2026}
Qin Liu.
``\$λ\_A\$: A Typed Lambda Calculus for LLM Agent Composition.''
arXiv:2604.11767, 2026.

\bibitem{llmbda2026}
Zac Garby et al.
``The LLMbda Calculus: AI Agents, Conversations, and Information Flow.''
arXiv:2602.20064, 2026.

\bibitem{resourcecontracts2026}
Qing Ye and Jing Tan.
``Agent Contracts: A Formal Framework for Resource-Bounded Autonomous AI Systems.''
arXiv:2601.08815, 2026.

\bibitem{skillverify2026}
Alfredo Metere.
``Methods for Formal Verification of Agent Skills: Three Layers Toward a Mechanically Checkable Capability-Containment Proof.''
arXiv:2605.23951, 2026.

\bibitem{toolcalculus2026}
Andreas Schlapbach.
``Formal Semantics for Agentic Tool Protocols: A Process Calculus Approach.''
arXiv:2603.24747, 2026.

\bibitem{uacp2026}
Arnab Mallick and Indraveni Chebolu.
``μACP: A Formal Calculus for Expressive, Resource-Constrained Agent Communication.''
arXiv:2601.00219, 2026.

\end{thebibliography}
